Zunächst konvertieren wir 1 kg in GeV. In natürlichen Einheiten gilt $\hbar = c = 1$, was bedeutet, dass Energie, Masse und Impuls dieselben Einheiten haben. Die Umrechnung von kg in GeV erfolgt über die Ruheenergie eines Kilogramms:

\[
1 \, \text{kg} = \frac{1 \, \text{kg} \cdot c^2}{1.60218 \times 10^{-10} \, \text{J/GeV}}
\]

Da $c = 3 \times 10^8 \, \text{m/s}$ und $1 \, \text{J} = 1 \, \text{kg} \cdot \text{m}^2/\text{s}^2$, erhalten wir:

\[
1 \, \text{kg} = \frac{(3 \times 10^8 \, \text{m/s})^2}{1.60218 \times 10^{-10} \, \text{J/GeV}} \approx 5.61 \times 10^{26} \, \text{GeV}
\]

Nun konvertieren wir 1 s in $\mathrm{GeV}^{-1}$. Da $c = 1$, ist $1 \, \text{s} = 1 \, \text{m}$ in natürlichen Einheiten. Die Umrechnung erfolgt über die Beziehung:

\[
1 \, \text{s} = \frac{1}{1.519 \times 10^{24} \, \text{GeV}}
\]

Jetzt berechnen wir das Gravitationskonstante $G_N$ in natürlichen Einheiten. In SI-Einheiten hat $G_N$ die Dimension $\mathrm{m}^3 \, \mathrm{kg}^{-1} \, \mathrm{s}^{-2}$. Um dies in natürlichen Einheiten umzurechnen, verwenden wir:

\[
G_N = 6.67 \times 10^{-11} \, \mathrm{m}^3 \, \mathrm{kg}^{-1} \, \mathrm{s}^{-2} = 6.67 \times 10^{-11} \left( \frac{1}{1.519 \times 10^{24} \, \text{GeV}} \right)^2 \left( \frac{1}{5.61 \times 10^{26} \, \text{GeV}} \right)
\]

Dies ergibt:

\[
G_N \approx 6.71 \times 10^{-39} \, \text{GeV}^{-2}
\]

Schließlich berechnen wir die Planck-Masse $M_{Pl}$:

\[
M_{Pl} = G_N^{-1/2} \approx (6.71 \times 10^{-39} \, \text{GeV}^{-2})^{-1/2} \approx 1.22 \times 10^{19} \, \text{GeV}
\]